\documentclass[tikz,border=8pt]{standalone}
\usepackage{ctex}
\usepackage{amsmath}
\usetikzlibrary{calc, arrows.meta, patterns}

\begin{document}
% 线面垂直判定定理示意图
% 直线 l ⊥ 平面内两条相交直线 m, n → l ⊥ 平面 α
\begin{tikzpicture}[
  x={(1cm,0cm)},
  y={(0.38cm,0.28cm)},
  z={(0cm,0.9cm)},
  thick,
]

% 平面参数
\def\px{3.0}  % 平面x方向宽
\def\py{2.6}  % 平面y方向深

% 平面底角
\coordinate (A0) at (0,0,0);
\coordinate (B0) at (\px,0,0);
\coordinate (C0) at (\px,\py,0);
\coordinate (D0) at (0,\py,0);

% 平面原点（直线交点）
\coordinate (O) at (1.5,1.3,0);

% 平面内两相交直线 m, n 的端点
\coordinate (Mm) at (0.2,0.2,0);  % m 的一端
\coordinate (Mp) at (2.8,2.4,0);  % m 的另一端
\coordinate (Nm) at (2.6,0.2,0);  % n 的一端
\coordinate (Np) at (0.3,2.4,0);  % n 的另一端

% 直线 l 的顶点（竖直向上）
\coordinate (Ltop) at (1.5,1.3,2.5);

% 1. 画平面（填色）
\fill[blue!10, draw=blue!40, fill opacity=0.7]
  (A0) -- (B0) -- (C0) -- (D0) -- cycle;
\node[blue!60, font=\small] at ($(C0) + (0.3, 0.2)$) {$\alpha$};

% 2. 画平面内两相交直线 m, n
\draw[orange!80!black, thick] (Mm) -- (Mp);
\node[orange!80!black, font=\small, above right] at (Mp) {$m$};

\draw[green!60!black, thick] (Nm) -- (Np);
\node[green!60!black, font=\small, above left] at (Np) {$n$};

% 交点 O 标记
\fill[black] (O) circle [radius=1.5pt];
\node[font=\small, below right] at ($(O)+(0.05,-0.05)$) {$O$};

% 3. 画直线 l（垂直于平面）
% 平面以下虚线（视觉效果）
\draw[red, thick] (O) -- (Ltop);
\node[red, font=\small, above] at (Ltop) {$l$};

% 4. 直角符号：l ⊥ m
% 在 m 方向取单位向量，在 l 方向取单位向量，画直角框
% m 方向：(2.8-0.2, 2.4-0.2) = (2.6, 2.2), 归一化
\pgfmathsetmacro{\mux}{(2.8-0.2)/sqrt((2.8-0.2)^2+(2.4-0.2)^2)}
\pgfmathsetmacro{\muy}{(2.4-0.2)/sqrt((2.8-0.2)^2+(2.4-0.2)^2)}
\pgfmathsetmacro{\bsz}{0.18}  % 直角符号大小

% l ⊥ m 直角符号（在交点 O 处）
\draw[gray, thin]
  ($(O) + \bsz*(\mux, \muy, 0)$)
  -- ($(O) + \bsz*(\mux, \muy, \bsz)$)
  -- ($(O) + \bsz*(0, 0, \bsz)$);

% 5. 直角符号：l ⊥ n
\pgfmathsetmacro{\nux}{(0.3-2.6)/sqrt((0.3-2.6)^2+(2.4-0.2)^2)}
\pgfmathsetmacro{\nuy}{(2.4-0.2)/sqrt((0.3-2.6)^2+(2.4-0.2)^2)}

\draw[gray, thin]
  ($(O) + \bsz*(\nux, \nuy, 0)$)
  -- ($(O) + \bsz*(\nux, \nuy, \bsz)$)
  -- ($(O) + \bsz*(0, 0, \bsz)$);

% 6. 底部文字说明
\node[font=\small, align=center] at (1.5, -0.7, 0)
  {$l \perp m$，$l \perp n$，$m \cap n = O$};
\node[font=\small, align=center] at (1.5, -1.1, 0)
  {$\Rightarrow l \perp \alpha$（线面垂直判定定理）};

\end{tikzpicture}
\end{document}
